\section{The surviving path regime for a named polynomial subclass}
\label{sec:surviving-path-polynomial}

Set
\begin{equation}
 \alpha_n^{\rm surv}:=n^{-2},
 \qquad
 \frac{\nu_n}{\sqrt{\alpha_n^{\rm surv}}}
 =\Theta(n^2).
 \label{eq:surviving-path-regime}
\end{equation}
The task and budget remain exactly
\cref{eq:path-exact-output-task,eq:concrete-rowrec-budget}.  This section
does not prove a lower bound for the full budgeted row-recurrence framework.
It closes only the following named polynomial escape routes after the complete
$n$-coordinate path system has been exposed.

\paragraph{The named subclass.}
Grant an execution any fixed positive diagonal $M$ and put
\begin{equation}
 T:=M^{-1/2}Q_nM^{-1/2},
 \qquad h:=M^{-1/2}b_n.
 \label{eq:diagonal-krylov-transform}
\end{equation}
Any actual graph-dependent work used to choose or encode $M$ remains charged,
although the cyclicity argument below holds even if $M$ is supplied.
For $r\geq0$, define
$\mathsf{DiagSpecPoly}(r)$ as follows.  The response part may select a set
$J$ of at most $r$ eigenpairs $(\lambda_j,u_j)$ of $T$ and return the exact
spectral correction
\begin{equation}
 z_J:=\sum_{j\in J}\lambda_j^{-1}
       \inner{u_j}{h}u_j.
 \label{eq:spectral-atom-correction}
\end{equation}
The remaining output in the $M^{1/2}$ coordinates must have the form
\begin{equation}
 z_J+p(T)(I-\Pi_J)h,
 \qquad
 \Pi_J:=\sum_{j\in J}u_ju_j^\top,
 \label{eq:diag-spec-poly-output}
\end{equation}
where $p(T)(I-\Pi_J)h$ is generated by full-vector applications of $T$ and
scalar linear combinations.  Every $T$ application includes one charged
full-vector application of $Q_n$ plus diagonal actions charged to recurrence
when $M$ is a supplied fixed map and to response when $M$ is a
graph-dependent preconditioner.  The emitted solution in the original
coordinates is $M^{-1/2}$ times
\cref{eq:diag-spec-poly-output}, with that final diagonal arithmetic and its
output cells charged separately.

This named atom convention charges at least one response operation to
construct or identify each spectral atom and at least one for every
application.  No dense atom is a unit-cost object: every additional
construction and application operation is charged to $C_{\rm resp}$, every
explicit encoded atom and scratch cell to $M_{\rm pers}$ or $M_{\rm tmp}$,
every cell used to maintain a complete current correction to $C_{\rm mat}$,
and every required final cell to $C_{\rm emit}$.  Thus
$r\leq C_{\rm resp}\leq B_{\rm resp}$ holds only by this named atom-cost
convention.  It is \emph{not} an inference that an arbitrary response of
charge $B_{\rm resp}$ has rank at most $B_{\rm resp}$.

The case $r=0$ includes exact polynomial and Chebyshev evaluation, scalar
momentum and restart polynomials, polynomial preconditioners when every
underlying $T$ application is counted, and standard Krylov/PCG with the fixed
positive diagonal $M$.  The subclass deliberately excludes rational or
nonlinear response, adaptive nonspectral maps, elimination, and supported
frontier recurrences.

\begin{theorem}[Endpoint cyclicity blocks the named full-vector escapes;
\ProvedStatus]
\label{thm:surviving-path-polynomial-obstruction}
Let $n\geq16$ and use exact algebraic-cell arithmetic.  For every positive
diagonal $M$, the vector $h$ in \cref{eq:diagonal-krylov-transform} is cyclic
for $T$: its relative minimal polynomial has degree $n$.  Consequently:
\begin{enumerate}[label=(\roman*)]
\item an exact $r=0$ polynomial, Chebyshev, scalar-momentum, polynomially
      preconditioned, or fixed-positive-diagonal PCG execution needs a
      polynomial of degree at least $n-1$;
\item an exact $\mathsf{DiagSpecPoly}(r)$ execution needs degree at least
      $n-r-1$ after its spectral correction;
\item under the named atom convention, the defining full-vector application
      rule, and the budget $\mathcal B_n$,
      $r\leq\lfloor\sqrt{\nu_n}\rfloor$, so the execution needs
      $\Omega(n)$ full-vector $T$ applications.  Since each such application
      costs $\Theta(\nu_n)$ recurrence work on the exposed path,
      \begin{equation}
       C_{\rm rec}=\Omega(n\nu_n)
       =\Omega\!\left(
          \frac{\nu_n}{\sqrt{\alpha_n^{\rm surv}}}
        \right).
       \label{eq:surviving-path-polynomial-work}
      \end{equation}
\end{enumerate}
The first two conclusions lower-bound polynomial degree and the number of
$Q_n/T$ applications, not total work.  The last conclusion uses the
subclass's explicit full-vector path-application rule to convert them to
recurrence work.  Degree alone does not give this work bound for a supported
implementation whose active prefix grows over time.
\end{theorem}

\begin{proof}
The matrix $T$ is symmetric tridiagonal and every subdiagonal entry is
nonzero.  For a fixed eigenvalue, its first coordinate determines an
eigenvector recursively.  If that coordinate were zero, the first row and
then every successive row would force the whole eigenvector to vanish.
Thus every eigenspace is one-dimensional, all $n$ eigenvalues are distinct,
and every eigenvector has nonzero first coordinate.  Because
$b_n$ is a nonzero multiple of the first coordinate vector, so is $h$ after
diagonal scaling.  It therefore has nonzero projection on every eigenvector
of $T$ and is cyclic.

For $r=0$, exactness would give a polynomial $p$ with
$p(\lambda_j)=1/\lambda_j$ at all $n$ distinct eigenvalues.  If
$\deg p\leq n-2$, then $q(t):=tp(t)-1$ has degree at most $n-1$ and $n$
distinct roots, yet $q(0)=-1$, a contradiction.  Hence
$\deg p\geq n-1$.  After correcting $r$ eigencomponents in
\cref{eq:spectral-atom-correction}, the same argument on the remaining
$n-r$ eigenvalues gives $\deg p\geq n-r-1$.

By definition, one full-vector $T$ application can raise polynomial degree by
at most one.  Under \cref{eq:concrete-rowrec-budget} and the named atom charge,
\[
 r\leq C_{\rm resp}\leq\lfloor\sqrt{\nu_n}\rfloor=o(n).
\]
For $n\geq16$, $n-r-1=\Omega(n)$.  Independently of how the diagonal actions
are classified, every full $T$ application contains a full $Q_n$ application
and hence $\Theta(\nu_n)$ recurrence arithmetic on the stored exposed path.
Multiplying the application lower bound by this charged cost proves
\cref{eq:surviving-path-polynomial-work}; it does not add repeated adjacency
exposure to $C_{\rm adj}$.
\end{proof}

\paragraph{Eleven-coordinate ledger.}
When a budget-admissible $\mathsf{DiagSpecPoly}(r)$ execution is coupled to
the same chosen sequential exposure trace, its ledger constraints are
\begin{equation}
 \mathcal C_{\mathcal A}=\bigl(
 \Theta(\nu_n),n,1,0,C_{\rm ctl},\Omega(n\nu_n),C_{\rm resp},
 M_{\rm pers},M_{\rm tmp},C_{\rm mat},\Theta(n)\bigr),
 \label{eq:surviving-named-subclass-resource-constraints}
\end{equation}
where $C_{\rm ctl},C_{\rm mat}\geq0$,
$r\leq C_{\rm resp}\leq\lfloor\sqrt{\nu_n}\rfloor$, and
$M_{\rm pers},M_{\rm tmp}\geq0$ with
$M_{\rm pers}+M_{\rm tmp}\leq32\nu_n$.  The fixed values in this display are
properties of that chosen access/output implementation; only the recurrence
coordinate is lower-bounded by the cyclicity theorem after exposure.

At \cref{eq:surviving-path-regime}, the standard full-vector exact-CG
implementation uses no response atoms and is budget-admissible.  With the
same sequential exposure convention as \cref{eq:path-cg-resource-vector}, its
resource vector is
\begin{equation}
 \mathcal C_{\rm CG,surv}=\bigl(
 \Theta(\nu_n),n,1,0,\Theta(n),\Theta(n\nu_n),0,
 \Theta(\nu_n),\Theta(\nu_n),\Theta(n^2),\Theta(n)\bigr).
 \label{eq:surviving-path-cg-resource-vector}
\end{equation}
Here $R_{\rm adj}=n$ is the cost of this chosen exposure implementation, not
an adaptivity lower bound.  The complete-vector writes are those of the
standard materialized implementation.  The same-task fully charged
tridiagonal response still has the eleven-coordinate vector
\cref{eq:path-ldl-resource-vector} and $\Theta(\nu_n)$ total nonmemory work;
it is the comparison response arm and exceeds the row-recurrence budget
$B_{\rm resp}$.

The result is a tight $\Theta(n\nu_n)$ cap/obstruction for the named subclass,
not the desired theorem for every execution in
$\mathsf{RowRec}[\mathcal B_n]$.  In particular, it says nothing about an
arbitrary nonlinear or adaptive response map, a rational factor, an implicit
path solve, or a local algorithm that discovers and grows support.  The
degree proof by itself does not lower-bound the cumulative work of supported
path Krylov; that would require a separate sum-of-growing-support argument.
The cyclicity statement concerns a completely exposed fixed-support exact system
only.  It has no finite-precision consequence.
